\ifcorrige
\appendix
\part*{Réponses brèves et pistes}\addcontentsline{toc}{part}{Réponses brèves et pistes}
\chapter*{Comment utiliser les réponses}\addcontentsline{toc}{chapter}{Comment utiliser les réponses}
Les indications ci-dessous donnent les résultats finaux ou les idées décisives. Elles ne remplacent pas une rédaction complète. Pour les questions de construction, de justification ou de modélisation, plusieurs réponses peuvent être valides si les hypothèses et les contrôles sont clairement exposés.
\chapter*{Chapitre 1 -- Les nombres réels et leurs opérations}
\begin{answerbox}
\textbf{Exercice 1.1 -- Les grandes familles de nombres.} $-7\in\Z$; $18/6=3\in\N$; $0{,}125=1/8\in\Q$; $\sqrt{49}=7\in\N$; $\sqrt3\in\R\setminus\Q$; $\pi-3\in\R\setminus\Q$. Exemples : $\sqrt2+(-\sqrt2)=0$ et $\sqrt2+\sqrt3$ irrationnel. $A\cap B=]1,3[$, $A\cup B=[-2,5]$, $A\setminus B=[-2,1]$.\par
\textbf{Exercice 1.2 -- Priorité des opérations et signes.} Premier calcul : $-9+4\cdot9/6=-3$. $(-2)^4=16$, $-2^4=-16$, $(-2)^{-2}=1/4$, $-2^{-2}=-1/4$. Ici $ab^2<0$, donc $-ab^2>0$, puis division par $c<0$ : $E<0$.\par
\textbf{Exercice 1.3 -- Fractions et nombres décimaux.} $5/12-7/18+1/9=5/36$. $0{,}375=3/8$; $1{,}\overline{27}=14/11$; $0{,}1\overline6=1/6$. Pour $8$ portions : $(3/4)(8/5)=6/5=1{,}2$ tasse.\par
\textbf{Exercice 1.4 -- Puissances, notation scientifique et racines.} $2^{12-5-2}=2^5=32$; $3^{5-2-3}=1$. Produit : $1{,}5\times10^4$. Radical : $6\sqrt5-4\sqrt5+3\sqrt5=5\sqrt5$. $\sqrt{x^2}$ est la racine carrée non négative, donc $|x|$.\par
\textbf{Exercice 1.5 -- La droite réelle comme modèle unificateur.} $|x-2|\le3\iff x\in[-1,5]$. À moins de $0{,}4$ de $-1{,}5$ : $|x+1{,}5|<0{,}4$, soit $x\in]-1{,}9,-1{,}1[$. Écart absolu $0{,}5$; relatif $0{,}5/5{,}3\approx9{,}43\%$.\par
\textbf{Exercice 1.6 -- Cohérence des règles de calcul.} Les conventions sont imposées par la cohérence du quotient des puissances. Les deux premières égalités sont seulement parfois vraies; la troisième est vraie sur les réels si $a,b\ge0$. Par exemple $(1+1)/(1+1)=1$ tandis que $1/1+1/1=2$.\par
\textbf{Exercice 1.7 -- Mesurer, estimer et contrôler un résultat.} Le quotient est proche de $20\cdot4/0{,}5=160$, donc entre $150$ et $170$; valeur $\approx156{,}45$. $v$ s'exprime en km/h et vaut $75$. $37\%<50\%$, donc le résultat doit être $<240$; déplacer la virgule : $0{,}37\cdot480=177{,}6$.\par
\medskip\textbf{Synthèse.} $\sqrt2$ reste irrationnel; l'écriture décimale est tronquée. Montant final : $84(1{,}15)(0{,}85)=82{,}11\,$\$. Aire centrale $100{,}44\ \mathrm{cm}^2$; encadrement $12{,}3\cdot8{,}0\le A\le12{,}5\cdot8{,}2$, soit $98{,}4\le A\le102{,}5$.
\end{answerbox}
\chapter*{Chapitre 2 -- Expressions algébriques, équations et inéquations}
\begin{answerbox}
\textbf{Exercice 2.1 -- Variables, monômes et polynômes.} Degrés $3,3,0$; degré total $3$; $P(2,-1)=-12-10+7=-15$. Réduction : $5x^2-8x+2$. Aire $2x^2+5x-3$, périmètre $6x+4$; il faut $x>1/2$.\par
\textbf{Exercice 2.2 -- Factorisation et fractions algébriques.} $6x(x-2)(x+2)$; $(x-5)^2$; $(2x-3)(2x+3)$. Fraction : $(x+3)(x-3)/[(x+3)(x-2)]=(x-3)/(x-2)$ avec $x\neq-3,2$. $f$ et $g$ ont la même formule simplifiée mais des domaines différents; elles ne sont pas égales sur $\R$.\par
\textbf{Exercice 2.3 -- Équations du premier degré.} $x=38/5$. Deuxième : $4(x-2)-3(2x+1)=60$, donc $x=-71/2$. Égalité des forfaits : $18+0{,}08n=27+0{,}03n$, donc $n=180$.\par
\textbf{Exercice 2.4 -- Inéquations et tableaux de signes.} $-6x+3>5x+7\iff-11x>4\iff x<-4/11$. Deuxième : $x\in[-1/2,4]$. Troisième : $x\in]-\infty,-2[\cup]3,+\infty[$.\par
\textbf{Exercice 2.5 -- Équations avec racines ou valeurs absolues.} Domaine $x\ge0$; après élévation au carré, $x^2-2x-3=0$, solutions candidates $3,-1$, seule $3$ est valide. $|3x-5|=7$ donne $x=4$ ou $x=-2/3$. $|x+1|<4\iff-5<x<3$; $|2x-3|\ge5\iff x\le-1$ ou $x\ge4$.\par
\textbf{Exercice 2.6 -- Équivalences et transformations à contrôler.} Ajouter la même quantité est toujours réversible; multiplier par $x$ ne l'est que si $x\neq0$; le carré et la racine exigent des conditions. $x^2-x-6=0$ donne $x=3,-2$, seule $3$ convient. L'équivalence correcte est $x^2=9\iff x=\pm3$.\par
\textbf{Exercice 2.7 -- Équivalence, implication et ensemble des solutions.} $S_1=\{-2,2\}$, $S_2=\{2\}$; $x=2\Rightarrow x^2=4$, mais pas la réciproque. Par exemple divisible par $12$ est suffisant; divisible par $2$ est nécessaire. $2x+5=11\iff2x=6\iff x=3$.\par
\textbf{Exercice 2.8 -- Résoudre sans perdre d'information.} Domaine $x\neq2$; $x+1=3x-6$, donc $x=7/2$. Deuxième : il faut $x\ge1$; carré donne $x+1=x^2-2x+1$, donc $x(x-3)=0$, seule $x=3$. Procédure : domaine, transformations, solutions candidates, substitution dans l'équation initiale.\par
\medskip\textbf{Synthèse.} Racines $2$ et $3$. Seuil $30x=120+18x$, donc $x=10$; bénéfice pour $x>10$. Troisième : $(x-1)-2(x+2)\le0$ sur $x\neq-2$, soit $(-x-5)/(x+2)\le0$, solution $]-\infty,-5]\cup]-2,+\infty[$.
\end{answerbox}
\chapter*{Chapitre 3 -- Fonctions : définition, graphe et opérations}
\begin{answerbox}
\textbf{Exercice 3.1 -- La notion de fonction.} Fonctions : $y=2x-1$ et $y=|x|$. Les deux autres peuvent associer deux valeurs de $y$ à un même $x$. Domaine : $[-2,1[\cup]1,+\infty[$. Entrée $C$ en degrés Celsius; sortie $F$ en degrés Fahrenheit.\par
\textbf{Exercice 3.2 -- Graphe et lecture graphique.} Racines $1,3$; ordonnée à l'origine $3$; sommet $(2,-1)$; axe $x=2$. $f(x)=2$ correspond aux intersections avec la droite $y=2$; $f(x)>2$ aux portions au-dessus. Le contexte restreint le domaine au temps observé.\par
\textbf{Exercice 3.3 -- Transformations de graphes.} Réflexion verticale, étirement facteur $2$, translation de $3$ à droite et $1$ vers le haut. Points : $(a,b+4)$, $(a+2,b)$, $(a,-b)$, $(-a,b)$. Pour \(f(x)=|x-1|\), \(f(2x)=|2x-1|\) a son zéro en \(1/2\), tandis que \(2f(x)=2|x-1|\) a son zéro en \(1\). La première transformation comprime horizontalement; la seconde étire verticalement.\par
\textbf{Exercice 3.4 -- Opérations et composition.} $\dom(f+g)=\dom(fg)=[-1,2[\cup]2,+\infty[$; pour $f/g$, il faut en plus $g(x)\neq0$, toujours vrai sur son domaine, donc même domaine. $f\circ g=2x^2-1$; $g\circ f=(2x-3)^2+1$. Les facteurs $1{,}15$ et $0{,}8$ commutent, donc même prix final $0{,}92x$.\par
\textbf{Exercice 3.5 -- Fonction réciproque.} $f^{-1}(x)=(x+7)/3$. Pour $x^2$, le test horizontal échoue; sur $[0,+\infty[$, inverse $\sqrt{x}$. Rationnelle : $f^{-1}(x)=(3x+1)/(x-2)$; domaine de $f$ : $\R\setminus\{3\}$, image $\R\setminus\{2\}$.\par
\textbf{Exercice 3.6 -- Situation, formule, tableau et graphe.} $V(t)=120-8t$, domaine contextuel $0\le t\le15$. Deuxième : différences constantes $3$, donc $f(x)=3x+4$. Troisième : un réservoir de \(50\) L perd \(3\) L/min; \(t\) est en minutes et le modèle cesse à \(t=50/3\), lorsque le réservoir est vide.\par
\textbf{Exercice 3.7 -- Lire une fonction à partir de son graphe.} Construction variable. Un zéro est une intersection avec l'axe des $x$; l'ordonnée à l'origine est $f(0)$; local compare un voisinage, global tout le domaine. Les intersections avec $y=x$ sont les points fixes de $f$.\par
\textbf{Exercice 3.8 -- Analyser une fonction sous plusieurs représentations.} $f(x)=3-x$ si $x<2$, $x-1$ si $x\ge2$; sommet $(2,1)$. Exemples de modèles : $3-\frac12x^2$ et un polynôme de degré supérieur passant par les trois points. Développée : ordonnée à l'origine; factorisée : racines; canonique : sommet.\par
\medskip\textbf{Synthèse.} Exemple : $f(x)=|x-1|$; \(f(x)=1-x\) si \(x<1\), et \(f(x)=x-1\) si \(x\ge1\); les deux demi-droites se rejoignent en \((1,0)\) avec un coin. L'égalité exige même domaine, codomaine et mêmes valeurs partout. Si $g(x)=rx$ convertit et $h(y)=y+c$ ajoute les frais, la composition est $h\circ g(x)=rx+c$.
\end{answerbox}
\chapter*{Chapitre 4 -- Modèles linéaires et quadratiques}
\begin{answerbox}
\textbf{Exercice 4.1 -- Le modèle affine.} Pente $2$, équation $y=2x+1$. Taxi $C(d)=4{,}25+1{,}85d$; $C(12)=26{,}45$; $d=25{,}75/1{,}85\approx13{,}92$. Intersection $x=6$; $f>g$ pour $x>6$.\par
\textbf{Exercice 4.2 -- Le modèle quadratique.} $f(x)=2(x-2)^2-3$; minimum $-3$ en $x=2$. $A(x)=-x^2+12x$, domaine $[0,12]$, maximum $36$ en $x=6$. Forme $a(x+2)(x-5)$; $-10a=-20$, donc $a=2$.\par
\textbf{Exercice 4.3 -- Zéros et formule quadratique.} $\Delta=121$, racines $3$ et $-2/3$; somme $7/3$, produit $-2$. Deuxième $\Delta=-24<0$, aucun zéro réel. Racine double si $36-4k=0$, donc $k=9$.\par
\textbf{Exercice 4.4 -- Modéliser et interpréter.} $h(0)=1{,}5$; sommet à $t=2$; $h(2)=21{,}1$. Critères : mécanisme attendu, simplicité, domaine d'utilisation, comportement futur. Un temps négatif peut être hors du domaine physique.\par
\textbf{Exercice 4.5 -- Voir la variation : différences premières et secondes.} Différences $5,7,9,11$, secondes $2,2,2$: modèle quadratique, ici $x^2+4x+2$. Calcul algébrique donne $f(x+2)-2f(x+1)+f(x)=2a$. Rapports constants suggèrent un modèle exponentiel.\par
\textbf{Exercice 4.6 -- Passer de données à un modèle.} $y=3x+3$, prédiction $27$. Deuxième $y=(x+1)^2$. Une structure systématique des résidus indique que le modèle ne capture pas un mécanisme ou une courbure.\par
\textbf{Exercice 4.7 -- Construire, comparer et valider un modèle.} Résidus $2$ et $-2$; mêmes erreurs absolues. Le modèle peut rester utile localement mais ne doit pas être extrapolé au-delà du domaine où les prédictions ont un sens. Démarche : variables/unités, ajustement, résidus, plausibilité, validation/extrapolation.\par
\medskip\textbf{Synthèse.} Coût $C(x)=15x+120$, donc $C(35)=645$. Pour $h=17$, $-5t^2+20t-15=0$, donc $t=1$ ou $3$; maximum $22$ à $t=2$; sol à $t=2+\sqrt{110}/5\approx4{,}10$. Jeux de données variables, mais différences premières constantes pour l'affine et secondes constantes pour le quadratique.
\end{answerbox}
\chapter*{Chapitre 5 -- Fonctions polynomiales et rationnelles}
\begin{answerbox}
\textbf{Exercice 5.1 -- Architecture des fonctions polynomiales et rationnelles.} Degré $4$, coefficient $-2$, fonction paire, $P(x)\to-\infty$ aux deux extrémités. $R=(x-2)(x+2)/[(x-3)(x+2)]$, domaine $\R\setminus\{-2,3\}$, simplification $(x-2)/(x-3)$ avec trou en $-2$. Numérateur : zéros; dénominateur : restrictions et asymptotes possibles.\par
\textbf{Exercice 5.2 -- Prévoir un graphe à partir de l'expression.} Racines $-2$ (touche) et $1$ (traverse avec aplatissement); degré $5$, coefficient dominant négatif : gauche $+\infty$, droite $-\infty$. Pour $g$: asymptote verticale $x=3$, horizontale $y=2$, zéro $x=-1/2$, ordonnée $-1/3$. Exemple $(x+1)^2(x-2)^3$.\par
\textbf{Exercice 5.3 -- Comportement local et comportement global.} Puissance impaire traverse; paire touche; la puissance $3$ traverse plus à plat. $1/(x-2)\to-\infty$, $+\infty$, puis $0$. Pour $|x|$ grand, $x^5$ domine; près de $0$, les autres termes peuvent être comparables ou dominants.\par
\textbf{Exercice 5.4 -- Du facteur au graphe.} Signe positif sur $]-\infty,-3[$, négatif sur $]-3,1[$ et $]1,4[$, positif sur $]4,+\infty[$; pas de changement en $1$. Division : $x^2+1=(x-1)(x+1)+2$, asymptote $y=x+1$. Pour $R$: domaine exclut $\pm2$, zéro $-1$, asymptotes verticales $x=\pm2$, horizontale $y=0$; signes par facteurs.\par
\medskip\textbf{Synthèse.} Pour $f$: domaine exclut $\pm2$, zéros $\pm1$, asymptotes verticales $\pm2$, horizontale $y=1$, fonction paire. Deuxième : $P(x)=a(x+2)^2(x-3)$; $P(0)=-12a=-12$, donc $a=1$. Troisième : réponses variées, par exemple $(x+1)/(x-1)$ et $(2x-3)/(2x+5)$ ont asymptote $y=1$.
\end{answerbox}
\chapter*{Chapitre 6 -- Fonctions exponentielles et logarithmiques}
\begin{answerbox}
\textbf{Exercice 6.1 -- Variation additive et variation multiplicative.} $u_1=108,u_5=140,u_{10}=180$; $v_1=108,v_5\approx146{,}93,v_{10}\approx215{,}89$. Différences constantes : affine; rapports constants : exponentielle. Facteur $0{,}88$, modèle $Q_n=Q_0(0{,}88)^n$.\par
\textbf{Exercice 6.2 -- Le logarithme comme fonction réciproque.} $\log_2 32=5$; $\log_{10}(0{,}001)=-3$; $3^4=81$. $x=\log_5 17=\ln17/\ln5\approx1{,}760$. Les coordonnées $(x,b^x)$ sont échangées par la réciproque.\par
\textbf{Exercice 6.3 -- Lois des logarithmes.} $3\log_bx+\frac12\log_by-2\log_bz$ pour $x,y,z>0$. $\ln(3x^2/\sqrt y)$. Par exemple $x=y=1$: $\ln2\neq0$.\par
\textbf{Exercice 6.4 -- Équations exponentielles et logarithmiques.} $x=(1+\log_3 7)/2$. Domaine $x>1$; $(x-1)(x+2)=12$, soit $x^2+x-14=0$, seule $(-1+\sqrt57)/2>1$ convient. $t=\ln(7{,}5)/0{,}4\approx5{,}04$.\par
\textbf{Exercice 6.5 -- Applications.} Valeur $5000(1{,}042)^8\approx6948{,}83\,$\$; doublement $\ln2/\ln1{,}042\approx16{,}84$ ans. Fraction $(1/2)^{15/6}\approx0{,}1768$. La concentration au pH $3$ est $100$ fois celle au pH $5$.\par
\textbf{Exercice 6.6 -- Interprétation d'un facteur multiplicatif.} Valeur initiale $240$, croissance $3\%$ par période. $k=\ln1{,}03\approx0{,}02956$. Facteur $0{,}85$ signifie diminution relative de $15\%$; la perte absolue dépend de la valeur courante.\par
\textbf{Exercice 6.7 -- Trois façons de voir une exponentielle.} La quantité double tous les $5$ temps; tableau $80,160,320,640$. $\ln y=\ln A+x\ln b$. Rapports $2$; $\ln y=\ln3+x\ln2$.\par
\textbf{Exercice 6.8 -- Modéliser une variation multiplicative.} $b=(1{,}5)^{1/4}\approx1{,}1067$. $k=\ln(80/30)/6\approx0{,}1635\ \mathrm h^{-1}$. Le terme exponentiel tend vers $0$; $T_a$ est l'équilibre, $T_0-T_a$ l'écart initial.\par
\medskip\textbf{Synthèse.} La dernière équation se résout numériquement : $1000(1{,}05)^t-1000(1+0{,}05t)=200$, environ $t\approx12{,}11$ ans. Culture : $400\cdot3^{t/7}=10^6$, donc $t=7\ln2500/\ln3\approx49{,}85$ h. Troisième : méthode $b=(y_2/y_1)^{1/(x_2-x_1)}$, puis $A=y_1/b^{x_1}$.
\end{answerbox}
\chapter*{Chapitre 7 -- Principes de dénombrement}
\begin{answerbox}
\textbf{Exercice 7.1 -- Le principe fondamental du dénombrement.} Repas complets $4\cdot6\cdot3=72$; sans dessert $4\cdot6=24$. Codes $26^2\cdot10^3=676000$. L'arbre a $3\cdot4=12$ feuilles.\par
\textbf{Exercice 7.2 -- Factorielle et permutations.} $8!/6!=8\cdot7=56$. Livres : $7!=5040$; bloc : $6!\cdot2=1440$. MATHEMATIQUE contient $11$ lettres avec M, A, T chacun répété deux fois : $11!/(2!2!2!)=4\,989\,600$.\par
\textbf{Exercice 7.3 -- Arrangements.} $12\cdot11\cdot10=1320$. $A_{10}^4=10\cdot9\cdot8\cdot7=5040$. Sans répétition $8\cdot7\cdot6\cdot5=1680$; avec répétition $8^4=4096$.\par
\textbf{Exercice 7.4 -- Combinaisons.} $\binom{13}{5}=1287$. Choisir les $k$ retenus équivaut à choisir les $n-k$ non retenus. Équipes : $\binom52\binom94=10\cdot126=1260$.\par
\textbf{Exercice 7.5 -- Critères de sélection d'une méthode.} Mot de passe : produit; classer tous : permutation; comité : combinaison; médailles : arrangement. L'arbre suit les trois questions indiquées. $\binom{10}{3}3!=10!/(7!3!)\cdot3!=10!/7!=A_{10}^3$.\par
\textbf{Exercice 7.6 -- Comptage par complément et inclusion-exclusion.} $2^8-1=255$. Au moins une : $23+18-9=32$; aucune : $8$. Multiples : $50+20-10=60$.\par
\textbf{Exercice 7.7 -- Voir les choix comme un arbre.} $8$ feuilles; exactement deux faces : $3$ chemins. Trajets : $2\cdot3\cdot2=12$. Un arbre force l'énumération des choix successifs et distingue chaque chemin terminal.\par
\textbf{Exercice 7.8 -- Pourquoi l'ordre change le comptage.} Podiums $A_{10}^3=720$; comités $\binom{10}{3}=120$, facteur $6$. Il y a $ABC,ACB,BAC,BCA,CAB,CBA$. Main : $\binom{52}{4}$; séquence : $A_{52}^4$.\par
\textbf{Exercice 7.9 -- Construire un comptage fiable.} Premier chiffre $5$ choix, puis $5\cdot4\cdot3$, total $300$. Par inclusion-exclusion : $3^6-3\cdot2^6+3=540$. Contrôles : comparer au total sans restriction; tester la formule pour une petite taille énumérable.\par
\medskip\textbf{Synthèse.} PROBABILITE a $11$ lettres, avec B et I répétés deux fois : $11!/(2!2!)=9\,979\,200$. Voyelle initiale : traiter les choix de A, O, E et I avec attention; total $3\cdot10!/(2!2!)+10!/2!=4\,536\,000$. Deuxième : $\binom82\binom73+\binom83\binom72=28\cdot35+56\cdot21=2156$. Troisième : choisir les positions des zéros sans la première : $\binom52$; choisir et ordonner $4$ chiffres distincts non nuls : $A_9^4$; total $10\cdot3024=30240$.
\end{answerbox}
\chapter*{Chapitre 8 -- Modèles probabilistes élémentaires}
\begin{answerbox}
\textbf{Exercice 8.1 -- Expérience, univers et événement.} Univers $\{1,\ldots,6\}^2$; $A=\{(2,6),(3,5),(4,4),(5,3),(6,2)\}$. $B=\{FFF,FFP,FPF,PFF\}$; $C=\{FFF\}$. Intersection : figure rouge; union : rouge ou figure; complément : non rouge.\par
\textbf{Exercice 8.2 -- Axiomes et équiprobabilité.} Les axiomes sont satisfaits. Sommes $10,11,12$ : $3+2+1=6$ cas sur $36$, soit $1/6$. Exemple : choisir « pluie » ou « pas de pluie » comme deux issues ne les rend pas chacune de probabilité $1/2$.\par
\textbf{Exercice 8.3 -- Probabilité conditionnelle.} Il y a $12$ figures, dont $4$ rois : $1/3$. $P(M\mid F)=6/18=1/3$; $P(F\mid M)=6/9=2/3$. Produit : $P(A\cap B)=P(B)P(A\mid B)$.\par
\textbf{Exercice 8.4 -- Indépendance.} $P(A)=1/2$, $P(B)=1/2$, $P(A\cap B)=1/4$, donc indépendants. Incompatibles donnent $P(A\cap B)=0\neq P(A)P(B)>0$. Exactement une face : $2(0{,}6)(0{,}4)=0{,}48$.\par
\textbf{Exercice 8.5 -- Épreuve de Bernoulli et modèle binomial.} $X\sim\mathrm{Bin}(12,0{,}7)$. $P(X=8)=\binom{12}{8}0{,}7^8 0{,}3^4\approx0{,}2311$. Au moins une : $1-0{,}8^5\approx0{,}67232$.\par
\textbf{Exercice 8.6 -- Espérance mathématique.} $E(X)=0+0{,}5+0{,}6=1{,}1$. Gain net : $0{,}1(16)+0{,}9(-4)=-2\,$\$. $np$ est la moyenne à long terme du nombre de réussites, pas nécessairement une valeur observable sur un essai.\par
\textbf{Exercice 8.7 -- Proportion, fréquence et probabilité.} Estimations $0{,}63$ et $0{,}617$; la seconde repose sur plus de données. La fréquence est une donnée empirique; la probabilité est un paramètre du modèle. La fréquence tend à se stabiliser, sa variabilité relative diminue.\par
\textbf{Exercice 8.8 -- Le conditionnement avec des effectifs concrets.} Tableau : positifs $24$ malades, $6$ non malades; négatifs $6$ malades, $164$ non malades. $P(M\mid+)=24/30=0{,}8$; $P(+\mid M)=24/30=0{,}8$ ici par coïncidence; $P(M)=30/200=0{,}15$. En général les dénominateurs diffèrent et les probabilités ne sont pas interchangeables.\par
\textbf{Exercice 8.9 -- Interpréter une information probabiliste.} Réduction absolue $1$ point de pourcentage; relative $50\%$. Il manque le niveau initial et souvent la taille d'échantillon. Questions : méthode d'échantillonnage, taille et marge d'erreur, formulation et population cible.\par
\medskip\textbf{Synthèse.} Sur $10\,000$ : $100$ malades, $95$ vrais positifs; $9900$ non malades, $495$ faux positifs; $P(M\mid+)=95/590\approx16{,}1\%$. Succès en trois essais : $1-(5/6)^3=91/216$. Nombre moyen de lancers : $1+5/6+(5/6)^2=91/36\approx2{,}528$. Prix équitable : $10(0{,}2)+3(0{,}5)=3{,}50\,$\$.
\end{answerbox}
\chapter*{Chapitre 9 -- Vecteurs dans \texorpdfstring{$\R^2$ et $\R^3$}{R2 et R3}}
\begin{answerbox}
\textbf{Exercice 9.1 -- Du point au vecteur.} $\vect{AB}=(6,-4)$ et $\vect{BA}=(-6,4)$. Exemples : de $(0,0)$ à $(3,2)$ et de $(1,-1)$ à $(4,1)$. Un point est une position; un vecteur est un déplacement libre.\par
\textbf{Exercice 9.2 -- Opérations sur les vecteurs.} $u+v=(-3,1)$, $u-v=(7,-7)$, $3u=(6,-9)$, $-2v=(10,-8)$. Diagonale $u+v=(4,3)$. Déplacement total $(3,-1)$.\par
\textbf{Exercice 9.3 -- Norme, distance et vecteur unitaire.} $\norm u=10$, unitaire $(3/5,-4/5)$. Distance $\sqrt{3^2+4^2+(-4)^2}=\sqrt{41}$. Vecteurs $(3,4)$ et $(-3,-4)$.\par
\textbf{Exercice 9.4 -- Produit scalaire et angle.} Produit $1$; $\cos\theta=1/(\sqrt5\sqrt{10})=1/\sqrt{50}$, donc $\theta\approx81{,}87^\circ$. $3k-12=0$, donc $k=4$. La dernière identité devient $c^2=a^2+b^2-2ab\cos\theta$.\par
\textbf{Exercice 9.5 -- Combinaisons linéaires.} Résoudre $a+2b=7$, $a-b=1$: $a=3$, $b=2$. Pour le second, $(a,b,a+b)=(1,2,3)$, donc oui avec $a=1,b=2$. Deux directions indépendantes permettent d'atteindre toute combinaison de coordonnées du plan.\par
\textbf{Exercice 9.6 -- Interprétation géométrique : un déplacement libre.} Images : $(-1,2)$, $(1,2)$, $(0,5)$. Oui, un vecteur est libre si direction, sens et norme coïncident. $P'=P+\vect{AB}$ en coordonnées.\par
\textbf{Exercice 9.7 -- Le produit scalaire comme longueur d'une ombre.} Composantes scalaires $4$ et $3$. $\operatorname{proj}_v u=(u\cdot v/\norm v^2)v=(11/5)(1,2)=(11/5,22/5)$. Travail $50\cdot8\cos60^\circ=200$ J; composante utile $25$ N.\par
\textbf{Exercice 9.8 -- Relier déplacements et coordonnées.} Position finale $(4,3,9)$. Milieu $((x_A+x_B)/2,(y_A+y_B)/2)$, et analogue en $\R^3$. $P=A+\frac23(B-A)=(5,4)$.\par
\medskip\textbf{Synthèse.} $AB=\sqrt{20}$, $AC=\sqrt{37}$, $BC=\sqrt{41}$; $\cos A=(4\cdot1+2\cdot6)/(\sqrt{20}\sqrt{37})=16/\sqrt{740}$; aire $|4\cdot6-2\cdot1|/2=11$. Projection parallèle : $(u\cdot v/\norm v^2)v=(12/8)(2,2)=(3,3)$; orthogonale $(2,-2)$. $\vect{AB}=(2,-1,4)$, $\vect{AC}=(4,-2,8)=2\vect{AB}$, donc alignés.
\end{answerbox}
\chapter*{Chapitre 10 -- Équations de la droite et du plan}
\begin{answerbox}
\textbf{Exercice 10.1 -- Droite dans le plan.} $(x,y)=(2,-1)+t(3,4)$, soit $x=2+3t$, $y=-1+4t$. Pour $AB=(6,-6)$, direction $(1,-1)$, cartésienne $x+y-1=0$. Pour $P$, $t=2$ par $x$, mais $y=7$ correspond bien à $-1+8=7$: oui.\par
\textbf{Exercice 10.2 -- Vecteur normal.} Normal $(3,-2)$; directeur $(2,3)$ ou $(-2,-3)$. Équation $2(x-4)-5(y-1)=0$, soit $2x-5y-3=0$. Produit $(a,b)\cdot(b,-a)=ab-ab=0$.\par
\textbf{Exercice 10.3 -- Positions relatives de deux droites.} Directions colinéaires; $(5,0)=(1,2)+2(2,-1)$, donc droites confondues. Intersection : $x=3,y=1$. Directions colinéaires : parallèles ou confondues selon un point commun; sinon sécantes dans le plan.\par
\textbf{Exercice 10.4 -- Plan dans \texorpdfstring{$\R^3$}{R3}.} $2(x-1)-3(y-2)+4(z+1)=0$, soit $2x-3y+4z+8=0$. Pour $P$, $6+0+4+8\neq0$, donc non. Paramétrique $(x,y,z)=s(1,0,2)+t(0,1,-1)$.\par
\textbf{Exercice 10.5 -- Intersection d'une droite et d'un plan.} Substitution : $1+2t+t+2-t=3+2t=6$, donc $t=3/2$, point $(4,3/2,1/2)$. Si $d\cdot n\neq0$, une intersection; si $=0$, droite parallèle ou contenue. Exemple variable, p. ex. plan $z=0$, droite $(t,1,0)$.\par
\textbf{Exercice 10.6 -- Point, direction et vecteur normal.} Paramétrique : un point et une direction non nulle; cartésienne : un point et un normal. Élimination : $5(x-1)+2(y-3)=0$, soit $5x+2y-11=0$. Point $(0,0,7/3)$ et directions orthogonales au normal, par exemple $(1,2,0)$ et $(0,3,1)$.\par
\textbf{Exercice 10.7 -- Choisir et exploiter une représentation.} La forme cartésienne permet une substitution directe; la paramétrique permet de résoudre pour le paramètre. Pour une intersection, substituer la paramétrique dans la cartésienne est souvent efficace. Stratégie variable mais doit justifier le choix.\par
\medskip\textbf{Synthèse.} $\vect{AB}=(2,1,-2)$, $\vect{AC}=(-1,2,-1)$; un normal est leur produit vectoriel $(3,4,5)$, donc $3x+4y+5z=13$. Pour l'intersection, soustraction donne $x-2y=-2$; poser $y=t$, alors $x=2t-2$, $z=5-3t$. Dernier : $d\cdot n=2-1-2=-1\neq0$, donc une intersection; substitution donne $-1-t=4$, $t=-5$.
\end{answerbox}
\chapter*{Chapitre 11 -- Matrices et transformations}
\begin{answerbox}
\textbf{Exercice 11.1 -- Définition et dimensions.} $A$ est $2\times3$ et $a_{23}=6$. Exemple variable. L'égalité matricielle est définie entrée par entrée, donc exige la même grille d'indices.\par
\textbf{Exercice 11.2 -- Addition et multiplication par un scalaire.} $2A-3B=\begin{pmatrix}2&-16\\9&-6\end{pmatrix}$. $1{,}05A$ augmente chaque prix de $5\%$. Les entrées correspondantes n'ont pas les mêmes indices et la somme n'est pas définie.\par
\textbf{Exercice 11.3 -- Produit matriciel.} $AB=\begin{pmatrix}8&-2\\3&-1\end{pmatrix}$; $BA=\begin{pmatrix}2&4\\3&5\end{pmatrix}$. Le nombre de colonnes de $A$ doit égaler le nombre de lignes de $B$. Un vecteur-ligne de quantités multiplié par un vecteur-colonne de prix donne le total.\par
\textbf{Exercice 11.4 -- Matrice identité et puissances.} Vérification directe. $A^2=\begin{pmatrix}1&2\\0&1\end{pmatrix}$, $A^3=\begin{pmatrix}1&3\\0&1\end{pmatrix}$, donc $A^n=\begin{pmatrix}1&n\\0&1\end{pmatrix}$. $T^n$ applique la même transformation $n$ fois.\par
\textbf{Exercice 11.5 -- Matrices comme transformations du plan.} Images : colonnes $(2,1)$ et $(-1,3)$. Triangle image : $(0,0),(2,1),(-1,3)$. Matrices : $\begin{pmatrix}1&0\\0&-1\end{pmatrix}$, $\begin{pmatrix}3&0\\0&1\end{pmatrix}$, $\begin{pmatrix}0&-1\\1&0\end{pmatrix}$.\par
\textbf{Exercice 11.6 -- Une matrice déplace une grille.} $A$ est un cisaillement horizontal. $B$ donne un rectangle $2\times1/2$, aire conservée $1$. Exemple d'écrasement : $\begin{pmatrix}1&0\\0&0\end{pmatrix}$; toutes les composantes verticales sont perdues.\par
\textbf{Exercice 11.7 -- Déterminant : aire, orientation et perte d'information.} Déterminants $10$ et $0$. Valeur absolue : facteur d'aire; signe négatif : inversion d'orientation. Exemple $\begin{pmatrix}-3&0\\0&2\end{pmatrix}$ : réflexion et étirements, facteur d'aire $6$.\par
\textbf{Exercice 11.8 -- Calculer, composer et interpréter une transformation.} $R=\begin{pmatrix}0&-1\\1&0\end{pmatrix}$; $RS=\begin{pmatrix}0&-1\\2&0\end{pmatrix}$, $SR=\begin{pmatrix}0&-2\\1&0\end{pmatrix}$. $RS(1,2)=(-2,2)$. Déterminants : $2=1\cdot2$.\par
\medskip\textbf{Synthèse.} Matrice $\begin{pmatrix}2&-1\\1&2\end{pmatrix}$, déterminant $5$, parallélogramme d'aire $5$, inverse $\frac15\begin{pmatrix}2&1\\-1&2\end{pmatrix}$. Les deux compositions ne coïncident pas; calcul matriciel. À gauche, combinaison des lignes; à droite, combinaison des colonnes.
\end{answerbox}
\chapter*{Chapitre 12 -- Systèmes d'équations linéaires}
\begin{answerbox}
\textbf{Exercice 12.1 -- Trois points de vue sur un système.} Forme $A\mathbf x=\mathbf b$ avec $A=\begin{pmatrix}2&1\\1&-1\end{pmatrix}$, $\mathbf b=(7,2)^T$; solution $(3,1)$. Aucune : droites parallèles distinctes; infinité : équations proportionnelles.\par
\textbf{Exercice 12.2 -- Élimination de variables.} Addition : $8x=24$, donc $x=3,y=2$. Multiplier la première par $4$ et la seconde par $3$, ou PPCM $12$. Toute solution satisfait les deux équations, donc aussi leur combinaison; l'opération est réversible.\par
\textbf{Exercice 12.3 -- Matrice augmentée et élimination de Gauss.} $L_2\leftarrow L_2-3L_1$ donne $[0,-7\mid-11]$, donc $y=11/7$, $x=13/7$. Opérations : échange, multiplication non nulle, ajout d'un multiple. $0=5$ : contradiction; $0=0$ : équation redondante et variable libre possible.\par
\textbf{Exercice 12.4 -- Déterminant en dimension deux.} Déterminant $12-12=0$; les droites sont parallèles distinctes ici, donc aucune solution. Aire $|\det A|$. $6k-6=0$, donc $k=1$.\par
\textbf{Exercice 12.5 -- Matrice inverse.} $\det A=1$, $A^{-1}=\begin{pmatrix}3&-1\\-5&2\end{pmatrix}$. Solution $(2,3)$. Une transformation qui écrase une dimension n'est pas injective et ne peut être défaite.\par
\textbf{Exercice 12.6 -- Règle de Cramer.} $D=11$, $D_x=22$, $D_y=33$, donc $(x,y)=(2,3)$. $D\neq0$ garantit l'unicité et permet la division. Cramer devient peu efficace quand la dimension augmente; l'élimination est préférable.\par
\textbf{Exercice 12.7 -- Interprétation géométrique d'un système.} Pentes différentes : une; parallèles distinctes : aucune; proportionnelles : infinité. Intersection exacte $(5/3,7/3)$. En $\R^3$: point unique, droite commune, plan commun, aucune intersection.\par
\textbf{Exercice 12.8 -- Pourquoi l'élimination conserve les solutions.} On récupère $E_2$ en ajoutant $3E_1$, donc équivalence. Multiplier par zéro détruit toute l'information de l'équation. Le troisième résultat dépend de l'exemple mais les solutions coïncident.\par
\textbf{Exercice 12.9 -- Résoudre et interpréter un système.} Élimination donne $(x,y,z)=(1,2,3)$. Construction variable, par exemple trois équations indépendantes satisfaites par ce point. Une variable de quantité peut exiger $x\ge0$; la solution algébrique est alors hors contexte.\par
\medskip\textbf{Synthèse.} Déterminant $1-k^2$. Si $k\neq\pm1$, solution $x=y=1/(1+k)$. Si $k=1$, infinité avec $x+y=1$; si $k=-1$, contradiction $x-y=1$ et $-x+y=1$. Deuxième : $(x,y)=(24,22)$. Troisième : réponses variées; une ligne nulle avec deux pivots pour trois inconnues donne un paramètre libre.
\end{answerbox}
\chapter*{Chapitre 13 -- Programmation linéaire à deux variables}
\begin{answerbox}
\textbf{Exercice 13.1 -- Variables, contraintes et fonction objectif.} Contraintes $4x+2y\le80$, $x\ge10$, $x,y\ge0$. Objectif $P=50x+30y$. Chaque coefficient est une ressource ou un profit par unité de produit.\par
\textbf{Exercice 13.2 -- Région réalisable.} Sommets $(0,0),(4,0),(2,4),(0,6)$. $(3,3)$ viole $2x+y\le8$; $(4,0)$ oui; $(1,5)$ oui.\par
\textbf{Exercice 13.3 -- Principe d'optimalité aux sommets.} Valeurs : $0,12,14,12$; maximum $14$ en $(2,4)$. Les lignes de niveau se déplacent parallèlement et touchent en dernier un sommet ou une arête. Optimum multiple si une ligne de niveau est parallèle à une arête optimale.\par
\textbf{Exercice 13.4 -- Configurations particulières du domaine réalisable.} Exemple $x\ge3$ et $x\le1$. Région $x,y\ge0$ : objectif $x+y$ non borné. Avec $x+y\le5$, la contrainte $x+y\le10$ est redondante.\par
\textbf{Exercice 13.5 -- Optimiser en faisant glisser une droite.} Les lignes ont équation $y=-4x+P$ et se déplacent dans la direction du gradient $(4,1)$. Sommets $(0,0),(3,0),(2,3),(0,4)$; valeurs de $4x+y$: $0,12,11,4$, maximum $12$ en $(3,0)$.\par
\textbf{Exercice 13.6 -- De la phrase à la région réalisable.} $x\ge2y$, $x+y<50$, $5\le y\le20$. Le point test, souvent $(0,0)$, indique le côté qui satisfait l'inégalité. Des quantités produites négatives n'ont pas de sens.\par
\textbf{Exercice 13.7 -- De la décision au modèle géométrique.} Modèle variable mais doit comporter deux variables non négatives, contraintes linéaires et coût. La validation vérifie toutes les contraintes, les unités, l'intégralité éventuelle et l'interprétation pratique.\par
\medskip\textbf{Synthèse.} Sommets $(0,0),(50,0),(36,28),(0,40)$; profits $0,2000,2840,2000$, maximum $(36,28)$. Pour l'optimum multiple, choisir un objectif parallèle à une contrainte active. Avec intégralité, il faut tester les points entiers voisins réalisables; le théorème des sommets continus ne suffit pas directement.
\end{answerbox}
\chapter*{Chapitre 14 -- Angles et cercle trigonométrique}
\begin{answerbox}
\textbf{Exercice 14.1 -- Degrés et radians.} $\pi/6$, $3\pi/4$, $-7\pi/6$, $5\pi/2$. Degrés : $22{,}5^\circ$, $150^\circ$, $-315^\circ$, $540/\pi\approx171{,}89^\circ$. La relation vient du demi-tour.\par
\textbf{Exercice 14.2 -- Le cercle unité.} Points $(1,0),(0,1),(-1,0),(0,-1)$. Pour $P$, $\cos\theta=-3/5$, $\sin\theta=4/5$, quadrant II. L'identité est exactement $x^2+y^2=1$.\par
\textbf{Exercice 14.3 -- Angles remarquables.} Valeurs usuelles. $\sin150^\circ=1/2$, $\cos225^\circ=-\sqrt2/2$, $\tan300^\circ=-\sqrt3$. Signes : I tous +; II sinus +; III tangente +; IV cosinus +.\par
\textbf{Exercice 14.4 -- Périodicité et symétries.} Angles équivalents : $5\pi/6$, $3\pi/4$, $\pi/3$. Valeurs $-\sqrt3/2$, $1/2$, $1/\sqrt3$. Les symétries du cercle donnent les parités; changer de demi-tour change les deux coordonnées de signe, donc leur quotient reste le même.\par
\textbf{Exercice 14.5 -- Le radian mesure naturellement une rotation.} $s=r\theta=6$ cm. $\theta=18/12=1{,}5$ rad $\approx85{,}94^\circ$. Un radian est le même angle; la longueur d'arc vaut le rayon et change avec celui-ci.\par
\textbf{Exercice 14.6 -- Pourquoi les radians simplifient les formules.} En degrés, $s=r\pi\theta/180$; en radians, aucun facteur de conversion. Petit arc $0{,}06$ unité. Les radians sont sans dimension et mesurent directement le rapport arc/rayon.\par
\textbf{Exercice 14.7 -- Mesurer et lire une rotation.} $7\pi/3=1$ tour plus $\pi/3$, soit $7/6$ tour. Arc $10\pi/3$; secteur $\frac12r^2\theta=20\pi/3$. La calculatrice interprète $\pi/6$ comme degrés; il faut le mode radians.\par
\medskip\textbf{Synthèse.} Arc $8\cdot3\pi/4=6\pi$ cm; aire $\frac12\cdot64\cdot3\pi/4=24\pi$ cm$^2$. Points $(1/2,\sqrt3/2)$ et $(-1/2,\sqrt3/2)$. Oui, les angles diffèrent de $-2\pi$; même point final mais parcours orienté différent.
\end{answerbox}
\chapter*{Chapitre 15 -- Trigonométrie du triangle}
\begin{answerbox}
\textbf{Exercice 15.1 -- Triangles rectangles.} Autre côté $12$; angles $\arcsin(5/13)\approx22{,}62^\circ$ et $67{,}38^\circ$. Hauteur $40\tan52^\circ\approx51{,}20$ m. Choix respectifs : sinus, cosinus, tangente.\par
\textbf{Exercice 15.2 -- Vue d'ensemble des triangles quelconques.} Loi des cosinus : $c^2=49+81-126(1/2)=67$, donc $c=\sqrt{67}$. $C=75^\circ$; $b=8\sin70^\circ/\sin35^\circ\approx13{,}10$, $c\approx13{,}48$. SSS/SAS : cosinus; ASA/AAS : sinus; SSA : sinus avec multiplicité possible.\par
\textbf{Exercice 15.3 -- Choisir la formule à partir de la figure.} Équations selon la classification. Pour SSA, $h=b\sin A=3$; comme $h<a<b$ est faux puisque $a=10>b=6$, un seul triangle. Les contrôles indiqués doivent tous être satisfaits.\par
\textbf{Exercice 15.4 -- Origine géométrique des lois du sinus et du cosinus.} La hauteur vaut $b\sin A=a\sin B$, d'où le rapport. La projection de $b$ sur $a$ vaut $b\cos C$ et le calcul de distance donne la formule. Pour $C=90^\circ$, $\cos C=0$.\par
\textbf{Exercice 15.5 -- Résoudre et valider un triangle.} $c=\sqrt{12^2+15^2-2(12)(15)\cos70^\circ}\approx15{,}68$; puis $A\approx45{,}98^\circ$ et $B\approx64{,}02^\circ$. Deuxième : $\sin B=10\sin40^\circ/7\approx0{,}9183$, donc $B\approx66{,}69^\circ$ ou $113{,}31^\circ$; les deux donnent un angle $C$ positif, donc deux triangles. Contrôles : somme $180^\circ$, ordre côtés/angles, substitution dans une loi.\par
\medskip\textbf{Synthèse.} $C=65^\circ$; $a=120\sin48^\circ/\sin65^\circ\approx98{,}39$, $b\approx121{,}88$; hauteur $b\sin48^\circ\approx90{,}57$ m. Deuxième : angles par cosinus, aire par $\frac12ab\sin C$ ou Héron; plus grand angle opposé à $10$. Exemples variables respectant les critères $a<h$, $a=h$, $h<a<b$, $a\ge b$.
\end{answerbox}
\chapter*{Chapitre 16 -- Fonctions trigonométriques et fonctions réciproques}
\begin{answerbox}
\textbf{Exercice 16.1 -- Graphes fondamentaux.} Sinus/cosinus : domaine $\R$, image $[-1,1]$, période $2\pi$; sinus impair, cosinus pair. Tangente : domaine exclut $\pi/2+k\pi$, image $\R$, période $\pi$, impaire. $\sin x=\cos x$ donne $x=\pi/4,5\pi/4$.\par
\textbf{Exercice 16.2 -- Amplitude, période, déphasage et ligne médiane.} Amplitude $3$, période $\pi$, déphasage $\pi/4$ à droite, médiane $y=-1$. Exemple $y=3\cos(\pi(x-2)/3)+4$. $A$ amplitude/réflexion, $B$ période, $C$ translation horizontale, $D$ translation verticale.\par
\textbf{Exercice 16.3 -- Fonctions trigonométriques réciproques.} $\pi/6$, $3\pi/4$, $\pi/4$. $\arcsin$ retourne l'angle principal dans $[-\pi/2,\pi/2]$, donc $\pi/6$. Domaines/images usuels : arcsin $[-1,1]\to[-\pi/2,\pi/2]$; arccos $[-1,1]\to[0,\pi]$; arctan $\R\to]-\pi/2,\pi/2[$.\par
\textbf{Exercice 16.4 -- Équations trigonométriques élémentaires.} $x=\pi/3,2\pi/3$; général $x=\pi/3+2k\pi$ ou $2\pi/3+2k\pi$. Deuxième : $x=\pi/4,3\pi/4,5\pi/4,7\pi/4$. Troisième : $2x=\pi/4+k\pi$, donc $x=\pi/8+k\pi/2$; dans l'intervalle $\pi/8,5\pi/8$.\par
\textbf{Exercice 16.5 -- La sinusoïde est l'ombre d'une rotation.} $P(\theta)=(4\cos\theta,4\sin\theta)$; ordonnée $4\sin\theta$. Les valeurs successives sont $0,1,0,-1,0$ à rayon unité. Les coordonnées du point sont précisément ces projections signées.\par
\textbf{Exercice 16.6 -- Effet des paramètres d'une sinusoïde.} Transformations respectives : référence, étirement vertical, compression horizontale, décalage. Pour $f$: max $6$, min $-2$, période $2\pi/3$; max lorsque $\cos3x=-1$, soit $x=\pi/3+2k\pi/3$. Troisième : amplitude $3$, médiane $2$, $B=1/2$, modèle $3\sin((x-\pi)/2)+2$.\par
\textbf{Exercice 16.7 -- Construire une sinusoïde et retrouver un angle.} Médiane, amplitude, période, quart de période, phase. Pour la fonction : période $6$, points $(0,-1),(1{,}5,1),(3,-1),(4{,}5,-3),(6,-1)$. Équation $\sin(\pi t/3)=1/2$, donc $t=0{,}5$ et $2{,}5$ sur $[0,6]$.\par
\medskip\textbf{Synthèse.} Amplitude $2$, médiane $1$, période $2\pi/3$, déphasage $\pi/6$. Zéros : cosinus $=-1/2$; $f=2$ donne cosinus $=1/2$. Deuxième : amplitude $4$, médiane $8$, période $8$, modèle $8+4\cos(\pi(t-1)/4)$. Troisième : $\sqrt2\sin(x+\pi/4)=1$, solutions $x=0,\pi/2$ dans l'intervalle.
\end{answerbox}
\chapter*{Chapitre 17 -- Mouvement circulaire et modèles sinusoïdaux}
\begin{answerbox}
\textbf{Exercice 17.1 -- Vitesse angulaire.} $18\cdot2\pi/60=3\pi/5$ rad/s. $T=2\pi/5$ s, $f=5/(2\pi)$ Hz. $v=r\omega$, donc la grande roue a une vitesse linéaire plus grande.\par
\textbf{Exercice 17.2 -- Coordonnées d'un mouvement circulaire.} $x=2+6\cos(2t+\pi/3)$, $y=-1+6\sin(2t+\pi/3)$. À $t=\pi/4$, angle $5\pi/6$, position $(2-3\sqrt3,2)$. L'identité trigonométrique donne l'équation du cercle.\par
\textbf{Exercice 17.3 -- Construire un modèle à partir de données.} Amplitude $20$, médiane $22$, $\omega=2\pi/90=\pi/45$. Bas au départ : $h=22-20\cos(\pi t/45)$. Médiane montante : $h=22+20\sin(\pi t/45)$.\par
\textbf{Exercice 17.4 -- Interpréter un modèle.} Médiane $12$, amplitude $7$, période $8$, déphasage $3$ à droite. Médiane lorsque $t=3+4k$; max $19$ à $t=5+8k$; min $5$ à $t=9+8k$. Domaine de trois cycles : intervalle de longueur $24$ adapté au début des observations.\par
\textbf{Exercice 17.5 -- Interprétation ondulatoire.} Amplitude $3$, nombre d'onde $2\pi/5$ donc $\lambda=5$; $\omega=4\pi$, donc $f=2$ et vitesse $v=\omega/k=10$. Amplitude influence l'intensité; fréquence la hauteur. Période mesure le temps d'un cycle, longueur d'onde la distance d'un cycle.\par
\textbf{Exercice 17.6 -- Du mouvement circulaire au modèle temporel.} $h(t)=D+r\sin(\omega t+\phi)$ ou forme cosinus. Max $D+r$, min $D-r$, période $2\pi/|\omega|$. Cosinus convient à un départ à un extremum; sinus à un passage par la médiane.\par
\textbf{Exercice 17.7 -- Construire un modèle périodique à partir d'observations.} Période $8$, amplitude $6$, médiane $9$, modèle $9+6\cos(\pi(t-2)/4)$. Deuxième $7+4\sin(\pi(t-1)/6)$. Une phase peut donner la bonne valeur initiale mais le mauvais signe de pente; vérifier si le phénomène monte ou descend.\par
\textbf{Exercice 17.8 -- Construire et valider un modèle périodique.} Médiane par moyenne des extrêmes, amplitude par demi-étendue, période par distance entre pics; phase par un repère. Le modèle a médiane $18$, amplitude $5$, période $24$, maximum $23$ à $t=4$, minimum $13$ à $t=16$. Causes : bruit, asymétrie, amplitude variable, superposition de fréquences, dérive.\par
\medskip\textbf{Synthèse.} Amplitude $2$, médiane $2{,}8$, demi-période $6{,}2$ h, donc période $12{,}4$ h et modèle $2{,}8+2\cos(2\pi(t-3)/12{,}4)$. À $t=6$, hauteur $\approx2{,}90$ m. Roue : $120$ tr/min $=2$ tr/s, $\omega=4\pi$ rad/s, $T=0{,}5$ s, $v=0{,}35\cdot4\pi\approx4{,}40$ m/s. Revoir le déphasage si le décalage est constant; la période si l'erreur s'accumule au fil des cycles.
\end{answerbox}
\fi
\backmatter
\chapter*{Conclusion}
Ce cahier doit être utilisé avec le manuel : l'objectif n'est pas seulement d'obtenir une réponse, mais de choisir une représentation, justifier la méthode et contrôler le résultat.

